\documentclass[10pt,a4paper]{amsart}
\usepackage[margin=19mm]{geometry}
\usepackage{amsmath,amssymb,mathtools}
\usepackage[scr=rsfs,cal=euler]{mathalfa}
\usepackage{xfrac}
\usepackage[unicode,hidelinks,bookmarks=false]{hyperref}
\newcommand{\ie}{i.e.\ }
\newcommand{\cf}{cf.\ }
\newcommand{\resp}{respectively\ }
\newcommand{\Subsection}{\subsection}
\providecommand{\nobreakdash}{\nobreak\hbox{-}}
\setlength{\emergencystretch}{2em}
\multlinegap0pt
\usepackage{ragged2e}
\usepackage{a4wide}
\usepackage{mathabx}
\usepackage{bbm}
\usepackage{mathrsfs}
\usepackage{verbatim}
\usepackage{upgreek}
\usepackage{relsize}
\usepackage{dsfont}
\newcommand{\eps}{\varepsilon}
\newcommand{\ov}{\overline}
\newcommand{\id}{\textnormal{id}}
\newcommand{\mc}{\mathcal}
\newcommand{\mb}{\mathbb}
\newcommand{\mbm}{\mathbbm}
\newcommand{\mrm}{\mathrm}
\newcommand{\mscr}{\mathscr}
\newcommand{\mf}{\mathfrak}
\newcommand{\msf}{\mathsf}
\newcommand{\I}{\mathbbm{1}}
\newcommand{\st}{\;\vline\;}
\newcommand{\vp}{\varphi}
\newcommand{\Lvp}{\Lambda_{\vp}}
\newcommand{\Lpsi}{\Lambda_{\psi}}
\newcommand{\hvp}{\widehat{\vp}}
\newcommand{\Lhvp}{\Lambda_{\hvp}}
\newcommand{\hpsi}{\widehat{\psi}}
\newcommand{\md}{\operatorname{d}\!}
\newcommand{\cst}{\ifmmode \mathrm{C}^* \else $\mathrm{C}^*$\fi}
\newcommand{\la}{\langle}
\newcommand{\ra}{\rangle}
\newcommand{\bbGamma}{{\mathpalette\makebbGamma\relax}}
\newcommand{\makebbGamma}[2]{%
  \raisebox{\depth}{\scalebox{1}[-1]{$\mathsurround=0pt#1\mathbb{L}$}}%
}
\newcommand{\NN}{\mathbb{N}}
\newcommand{\RR}{\mathbb{R}}
\newcommand{\KK}{\mathbb{K}}
\newcommand{\CC}{\mathbb{C}}
\newcommand{\ZZ}{\mathbb{Z}}
\newcommand{\GG}{\mathbb{G}}
\newcommand{\HH}{\mathbb{H}}
\newcommand{\TT}{\mathbb{T}}
\newcommand{\PP}{\mathbb{P}}
\newcommand{\DD}{\mathbb{D}}
\newcommand{\EE}{\mathbb{E}}
\newcommand{\QQ}{\mathbb{Q}}
\newcommand{\vv}{\mathrm{V}}
\newcommand{\Vv}{\mathds{V}}
\newcommand{\vV}{\text{\reflectbox{$\Vv$}}\:\!}
\newcommand{\VV}{{\mathds{V}\!\!\!\! \;\!\! \;\!\! \text{\reflectbox{$\Vv$}}\:\!}}
\newcommand{\ww}{\mathrm{W}}
\newcommand{\WW}{{\mathds{V}\!\!\text{\reflectbox{$\mathds{V}$}}}}
\newcommand{\Ww}{\mathds{W}}
\newcommand{\wW}{\text{\reflectbox{$\Ww$}}\:\!}
\newcommand{\is}[2]{\left\langle#1\,\vline\,#2\right\rangle}
\newcommand{\isma}[2]{\bigl\langle#1\,\vline\,#2\bigr\rangle}
\newcommand{\ismaa}[2]{\langle#1\,|\,#2\rangle}
\newcommand{\si}[2]{\bigl| #1\bigl\rangle\bigr\langle #2 \bigr|}
\newcommand{\sima}[2]{| #1\rangle\langle #2 |}
\newcommand{\wot}{\ifmmode \textsc{wot} \else \textsc{wot}\fi}
\newcommand{\sot}{\ifmmode \textsc{sot} \else \textsc{sot}\fi}
\newcommand{\sots}{\ifmmode \textsc{sot}^* \else \textsc{sot}$^*$\fi}
\newcommand{\ssot}{\ifmmode \sigma\textsc{-sot} \else $\sigma$-\textsc{sot }\fi}
\newcommand{\ssots}{\ifmmode \sigma\textsc{-sot}^* \else $\sigma$-\textsc{sot }$^*$\fi}
\newcommand{\swot}{\ifmmode \sigma\textsc{-wot} \else $\sigma$-\textsc{wot}\fi}
\newcommand{\Linfd}{\operatorname{L}^{\infty}(\wh\bbGamma)}
\newcommand{\Ljd}{\operatorname{L}^{1}(\wh\bbGamma)}
\newcommand{\CGd}{\mathrm{C}(\wh\bbGamma)}
\newcommand{\wh}{\widehat}
\newcommand{\wc}{\widecheck}
\newcommand{\wt}{\widetilde}
\newcommand{\whG}{\widehat{\GG}}
\newcommand{\whH}{\widehat{\HH}}
\newcommand{\LdG}{\operatorname{L}^{2}(\GG)}
\newcommand{\LdIrr}{\operatorname{L}^{2}(\IrrG)}
\newcommand{\oon}{\operatorname}
\newcommand{\lec}{\preccurlyeq}
\newcommand{\lecq}{\preccurlyeq_q}
\newcommand{\gec}{\succcurlyeq}
\newcommand{\gecq}{\mathbin{_q\!\gec}}
\renewcommand{\restriction}{\mathord{\upharpoonright}}
\newcommand{\rest}{\restriction}
\DeclareMathOperator{\lin}{span}
\DeclareMathOperator{\Irr}{Irr}
\DeclareMathOperator{\Pol}{Pol}
\DeclareMathOperator{\supp}{supp}
\DeclareMathOperator{\Tr}{Tr}
\DeclareMathOperator{\B}{B}
\DeclareMathOperator{\A}{A}
\DeclareMathOperator{\F}{F}
\DeclareMathOperator{\M}{M}
\DeclareMathOperator{\K}{K}
\DeclareMathOperator{\T}{T}
\DeclareMathOperator{\N}{N}
\DeclareMathOperator{\Proj}{Proj}
\DeclareMathOperator{\Dom}{Dom}
\DeclareMathOperator{\Graph}{Graph}
\DeclareMathOperator{\Mor}{Mor}
\DeclareMathOperator{\End}{End}
\DeclareMathOperator*{\esssup}{ess\,sup}
\DeclareMathOperator{\LL}{L}
\DeclareMathOperator{\HS}{HS}
\DeclareMathOperator{\Diag}{Diag}
\DeclareMathOperator{\Dec}{Dec}
\DeclareMathOperator{\Rep}{Rep}
\DeclareMathOperator{\SU}{SU}
\DeclareMathOperator{\CB}{CB}
\DeclareMathOperator{\Inn}{Inn}
\DeclareMathOperator{\Aut}{Aut}
\DeclareMathOperator{\Ad}{Ad}
\newtheorem{theorem}{Theorem}[section]
\newtheorem*{theoremb}{Theorem}
\newtheorem{proposition}[theorem]{Proposition}
\newtheorem{lemma}[theorem]{Lemma}
\newtheorem{corollary}[theorem]{Corollary}
\newtheorem{remark}[theorem]{Remark}
\newtheorem*{notation}{Notation}
\newtheorem{definition}[theorem]{Definition}
\newtheorem{example}[theorem]{Example}
\numberwithin{equation}{section}
\begin{document}
\title[Scaling automorphisms of compact quantum groups]{Scaling automorphisms of compact quantum groups}
\author{Pleinlaan 2; 1050 Brussels; Jacek Krajczok}
\date{}
\maketitle
\noindent\emph{Source and licence.} Scaling automorphisms of compact quantum groups. Version arXiv:2608.07089v1 (CC BY 4.0). This material is distributed under the Creative Commons Attribution 4.0 International licence, http://creativecommons.org/licenses/by/4.0/, as recorded in the arXiv OAI licence metadata for this exact version and shown on the abstract page https://arxiv.org/abs/2608.07089v1. Full rights notice: licenses/arxiv-2608.07089-CC-BY-4.0.txt.\par\smallskip
\noindent\textit{Editorial scope.} Counted unit: the original \texttt{proposition} environment \texttt{prop1} at source lines 250--252 together with its complete proof at lines 256--258; the delivered document keeps source lines 244--259 so that every referenced label and every citation resolves inside it. Native editable source is the paper's own concatenated TeX; the AMS wrapper is editorial formatting only. No publisher class, upstream script or unrelated artwork is redistributed.
\section{Proof of Theorem A}\label{sec:preliminaries}

We will freely use results from the theory of compact quantum groups \cite{NeshveyevTuset}, following mostly notational conventions from this book. \\

We will establish the following slightly stronger statement. It shows that it is the SOT-continuity of the implementation, rather than second countability of $\GG$, that is crucial.


\begin{proposition}\label{prop1}
Let $\GG$ be a compact quantum group. Assume that there is a group of unitary operators $(u_t)_{t\in \RR}$ in $\LL^{\infty}(\GG)$ which is continuous in SOT and $\tau_t=\oon{Ad}(u_t)|_{\LL^{\infty}(\GG)}$ for $t\in\RR$. Then $\GG$ is of Kac type.
\end{proposition}


Assuming this result, we first deduce Theorem A:


\begin{proof}[Proof of Theorem A]
We assume that each $\tau_t$ is inner and $\LL^2(\GG)$ is separable, hence the theorem of Kallman \cite{KallmanInner} applies, and there is a SOT-continuous group $(u_t)_{t\in \RR}$ of unitary operators in $\LL^{\infty}(\GG)$ which implement $(\tau_t)_{t\in\RR}$. Proposition \ref{prop1} ends the proof.
\end{proof}


\begin{thebibliography}{99}
\bibitem[KallmanInner]{KallmanInner} R.~R. Kallman. \newblock Groups of inner automorphisms of von {N}eumann algebras. \newblock {\em J. Functional Analysis}, 7:43--60, 1971.
\bibitem[NeshveyevTuset]{NeshveyevTuset} S.~Neshveyev and L.~Tuset. \newblock {\em Compact quantum groups and their representation categories}, volume~20 of {\em Cours Sp\'{e}cialis\'{e}s [Specialized Courses]}. \newblock Soci\'{e}t\'{e} Math\'{e}matique de France, Paris, 2013.
\end{thebibliography}
\end{document}
